\documentclass{ctexart}
\usepackage{amsmath}
\usepackage{amssymb}
\usepackage{amsfonts}
\usepackage{geometry}
\usepackage{fancyhdr}
\usepackage{amsthm}
\pagestyle{plain}
\geometry{a4paper, left=2.5cm, right=2.5cm, top=2.5cm, bottom=2.5cm}

\usepackage{hyperref}
\usepackage{cleveref}
\usepackage{amssymb}
\usepackage{mathtools}
\newtheorem{theorem}{定理}
\newtheorem{lemma}{引理}
\newtheorem{proposition}{命题}
\hypersetup{
    colorlinks=true,
}
\crefname{figure}{图}{图}
\crefname{table}{表}{表}
\crefname{equation}{式}{式}
\author{何铭源}
\date{\today}
\title{数学建模第四次作业}
\begin{document}
\maketitle
\section{学生录取}
\subsection{求概率$f_k$, $g_k$}
\subsubsection{$f_k$}
$f_k$表示前$k$名中第一名的学生在总$n$学生中排第一的概率,也就是说总排名第一,第二的同学都在前$k$名中.

因此,前$k$名中有总排名第二的同学的概率为$\frac{k}{n}$,总排名第一的同学在后$n-k$名中的概率为$\frac{k - 1}{n - 1}$.

因此有:
\begin{equation}
    g_k=\frac{k}{n}\cdot \frac{k - 1}{n - 1}
\end{equation}

\subsubsection{$g_k$}
$g_k$表示前$k$名中第一名的学生在总$n$学生中排第二的概率,也就是说总排名第二的同学在前$k$名中,总排名第一的同学在$n-k$名中.

因此,前$k$名中有总排名第二的同学的概率为$\frac{k}{n}$,总排名第一的同学在后$n-k$名中的概率为$\frac{n-k}{n-1}$.

因此有:
\begin{equation}
    g_k=\frac{k}{n}\cdot \frac{n-k}{n-1}
\end{equation}
\subsection{记$v_k$为不录取前$k$名学生后，采用最优策略可能录取到综合素质第二名的学生的概率的
    最大值。试写出$v_k$满足的递推关系。}
分情况讨论第$k + 1$学生的情况：
\begin{enumerate}
    \item 该学生在前$k + 1$名中排名第一：

          这种情况的概率为$\frac{1}{k + 1}$,此时我们需要比较$f_k$和$v_{k + 1}$的大小,最优的情况取最大值,因此该情况下的贡献为$\frac{1}{k + 1}\max(f_k, v_k)$.
    \item 该学生在前$k + 1$名中排名第二：

          这种情况的概率为$\frac{1}{k + 1}$,此时我们需要比较$g_k$和$v_{k + 1}$的大小,最优的情况取最大值,因此该情况下的贡献为$\frac{1}{k + 1}\max(g_k, v_k)$.

    \item 该学生在前$k + 1$名中排名第三及以后：

          这种情况的概率为$\frac{k - 1}{k + 1}$,此时我们只能采用之前的最优策略,因此该情况下的贡献为$\frac{k - 1}{k + 1}v_k$.
\end{enumerate}
综合以上讨论,最后的递推关系为:
\begin{equation}
    v_{k}=\frac{1}{k + 1}\max(f_k, v_{k + 1})+\frac{1}{k + 1}\max(g_k, v_{k + 1})+\frac{k - 1}{k + 1}v_{k + 1}
\end{equation}
讨论边际情况,$v_0 = \max(\frac{1}{n}, v_1)$,$v_n = 0$.

\begin{equation}
    v_k = \begin{dcases}
        \max(\frac{1}{n}, v_1),                                                                                   & k = 0     \\
        \frac{1}{k + 1}\max(f_k, v_{k + 1}) + \frac{1}{k + 1}\max(g_k, v_{k + 1}) + \frac{k - 1}{k + 1}v_{k + 1}, & 0 < k < n \\
        0,                                                                                                        & k = n
    \end{dcases}
\end{equation}

\subsection{求$v_k$,并给出相应的最优策略}
由招聘问题,我们需要设置一个观察期阈值,观察前$k$名学生,在$k$之后发现相对成绩第一名或第二名的学生时,就立马录取，并结束面试.

现在我们来确定这个阈值

我们需要找到$k$,使得超过$k$之后,$f_{k + 1} > v_{k + 1}$,且$g_{k + 1} > v_{k + 1}$,此时,招收满足条件的同学收益最大.

在超过阈值之后，$v_k$的递推公式退化为
\begin{equation}
    v_k = \frac{1}{k + 1}(f_{k + 1} + g_{k + 1}) + \frac{k - 1}{k + 1}v_{k + 1}
\end{equation}

因为$f_{k + 1} + g_{k + 1} = \frac{k + 1}{n}$,代入得、
\begin{equation}
    v_k = \frac{1}{n} + \frac{k - 1}{k + 1}v_{k + 1}
\end{equation}

这个公式当$f_{k + 1} > v_{k + 1}$且$g_{k + 1} > v_{k + 1}$时成立,即$k > \lfloor \frac{n}{2} \rfloor$成立.
\begin{equation}
    v_k = \begin{cases}
        \frac{\lfloor \frac{n}{2} \rfloor (n - \lfloor \frac{n}{2} \rfloor)}{n(n-1)}, & 0 \le k < \lfloor \frac{n}{2} \rfloor \\
        \frac{k(n-k)}{n(n-1)},                                                        & \lfloor \frac{n}{2} \rfloor \le k < n
    \end{cases}
\end{equation}
\section{第二题}
\subsection{试求出招聘方对第$r$名接受面试的求职者提出意向并被接受，而该求职者恰为能力最强的求职者的概率}
第$r$名求职者为能力最强的求职者概率在他在前$r$名中是能力最强的条件概率为$\frac{r}{n}$,接受的概率为$p$,因此总概率为:
\begin{equation}
    P_r = \frac{r}{n}\cdot p
\end{equation}
\subsection{记$v_r$为从招聘方面试第$r$名求职者起，采用最优策略以录用到能力最强的求职者的概率的最大值。试写出$v_r$满足的递推关系}
与上一题类似，在面试了$r - 1$名求职者后，分情况讨论第$r$名求职者的情况：
\begin{enumerate}
    \item 该求职者为前$r$名中能力最强的：

          这种情况的概率为$\frac{1}{r}$,此时我们需要比较$P_{r}$和$v_{r}$的大小,最优的情况取最大值,因此该情况下的贡献为$\frac{1}{r}\max(\frac{pr}{n} + (1 - p)v_{r + 1}, \quad v_{r + 1})$.
    \item 该求职者不是前$r$名中能力最强的：

          这种情况的概率为$\frac{r - 1}{r}$,此时我们只能采用之前的最优策略,因此该情况下的贡献为$\frac{r - 1}{r}v_{r + 1}$.
\end{enumerate}

综合以上讨论,最后的递推关系为:
\begin{equation}
    v_r = \frac{r-1}{r} v_{r+1} + \frac{1}{r} \max \left( p\frac{r}{n} + (1-p)v_{r+1}, \quad v_{r+1} \right)
\end{equation}
\subsection{求$v_r$的表达式与$\displaystyle{\lim_{n \to \infty}v_1}$，并给出招聘方最优策略的具体描述}
当提出录用的收益比不提出录用的收益大的时候，我们遇到满足条件的就可以录用。
\begin{equation}
    \begin{aligned}
        p\frac{r}{n} + (1-p)v_{r+1} & \ge v_{r+1}   \\
        p\frac{r}{n}                & \ge p v_{r+1} \\
        \frac{r}{n}                 & \ge v_{r+1}   \\
    \end{aligned}
\end{equation}
此时递推式满足：
\begin{equation}
    v_r = \frac{r-p}{r} v_{r+1} + \frac{p}{n}
\end{equation}
当$n$很大的时候，可以令$x = \frac{r}{n}$，记 $v_r \approx V(x)$,那么原来的递推式可以近似为微分方程：
\begin{equation}
    \frac{dV}{dx} = \frac{p(1 - V(x))}{x}
\end{equation}
解得：
\begin{equation}
    V'(x) - \frac{p}{x}V(x) = -p
\end{equation}

解得：
\begin{equation}
    V(x) = \frac{p}{p-1}x + C x^p
\end{equation}
再通过$x = 1$时，$V(1) = 0$，解得$C = \frac{1}{1-p}$.因此有：
\begin{equation}
    V(x) = \frac{p}{1-p}(x^p - x)
\end{equation}
那么我们要找的观察期个数就是$x \ge V(x)$的值，也就是$x^* = V(x^*)$.
\begin{equation}
    x^* = p^{\frac{1}{1-p}}
\end{equation}
\begin{equation}
    \lim_{n \to \infty} v_1 = x^* = p^{\frac{1}{1-p}}
\end{equation}
\end{document}